\question[20]
برنامه زیر را در نظر بگیرید و به سوالات پاسخ دهید.

\begin{latin}
    \begin{lstlisting}[language=C]
#include <stdio.h>
#include <pthread.h>
#include <unistd.h>

int counter = 0;

void *thread_func(void *arg) {
    int id = *(int *)arg;
    for (int i = 0; i < 3; i++) {
        counter++;
        printf("Thread %d: counter = %d\n", id, counter);
        sleep(1);
    }
    pthread_exit(NULL);
}

int main() {
    pthread_t t1, t2;
    int id1 = 1, id2 = 2;
    
    pthread_create(&t1, NULL, thread_func, &id1);
    pthread_create(&t2, NULL, thread_func, &id2);
    
    pthread_join(t1, NULL);
    pthread_join(t2, NULL);
    
    printf("Final counter = %d\n", counter);
    return 0;
}
\end{lstlisting}
\end{latin}

\begin{parts}
    \part
    اگر خط
    \lr{\texttt{pthread\_join(t1, NULL);}}
    را حذف کنیم، چه اتفاقی می‌افتد؟

    \part
    تفاوت اصلی بین ایجاد یک نخ جدید با \texttt{pthread\_create} و ایجاد یک فرآیند جدید با \texttt{fork} چیست؟
\end{parts}
